% Ad Quintum Fratrem 1.2 --- headnote
% ad-quintum-fratrem-01-02, between 27 October and 10 December 59 BC, written at Rome

\textit{Cicero to his brother Quintus, written at Rome between
27 October and 10 December 59 BC. Quintus is at the close of
his proconsulship of Asia (the great instructional letter
\textit{Q. fr.} 1.1 of late 60 / early 59 BC opens this same
governorship; this letter closes it). Statius, Quintus's
freedman whose manumission has been so much complained of in
the contemporary Atticus letters, has just arrived at Rome.
\textsection 1--3 take up Statius: his coming would have been
disastrous had it coincided with Quintus's, but coming alone
it has at least let the Roman gossip exhaust itself before
Quintus appears. Cicero is candid: even granting Statius's
fidelity, the appearance of so favoured a freedman has fed
every detractor of Quintus, and what was before merely
hostility to Quintus's strictness has, since the manumission,
not lacked talkers.}

\textit{The body of the letter \textsection 4--14 is the
brother-praetor's catalogue of Quintus's mistakes in the
conduct of his governorship --- written in the brother's
voice, not the consular's: candid, often sharp, always
loving. \textsection 4--5 on Cicero's having patiently soothed
all the Greek complainants from Hermippus of Dionysopolis to
Nymphon of Colophon, against the Zeuxis-of-Blaundus matter
where Quintus's strictness had run to attempted entrapment
(``two Mysians sewn up in a sack at Smyrna''). \textsection
6--9 on the angry letters: the Catienus letter (``the cross
he sets up for himself, from which I had previously dragged
him down''), the Fabius letter (``burn the kidnapper and his
chick alive''), all written for ``salt and humour'' but read
in cold ink as monstrous. The model of restraint is the
Cyrus of Xenophon and Agesilaus, ``those kings under whose
supreme command no one ever heard a sharper word.''
\textsection 10--11 on the Flavius--Fundanius case: Quintus's
edict ordering the heir Flavius not to diminish the estate
before Fundanius is paid is unjust as a matter of law and
will cost Cicero a friendship; Pompey and Caesar both
commended Flavius. \textsection 12 reads as the centre of the
piece: Cicero apologizes for his earlier angry letter on
Hermias, written under the spell of Diodotus, Lucullus's
freedman.}

\textit{\textsection 15--16 are the political coda. The
spectacular incident: the young M. Cato, attempting to indict
Gabinius for bribery, found the praetors inaccessible for
days; he mounted the Rostra and called Pompey \textit{a
private dictator}. He was within an inch of being killed.
``From this you can see the condition of the whole
commonwealth.'' Yet \textsection 16 is unexpectedly hopeful:
the friendly tribunes-elect, the well-intentioned consuls,
the keen praetors-elect (Domitius, Nigidius, Memmius,
Lentulus); the old \textit{manus bonorum} on fire with zeal.
The hope of \textsection 16 will collapse with Clodius's
tribunate of 58 BC; the next surviving letter (\textit{Att.}
3.1) is from the road into exile.}
